\sectiontheme{jsyellow}
\section{JSON Schema on the JavaScript Side}
\label{sec:json-schema}

Every JS/TS dto already has a field-plan tree the compiler walks to render
a form (\S\ref{sec:js-compiler}); \texttt{lib/formgen} walks that same
tree a second way, into a renderer-agnostic JSON Schema instead of a
bespoke component. There are two ways to get one out.

\subsection{Two ways to get a schema}
\textbf{Per-dto, standalone} --- \texttt{emi js:rjsf} takes a single
\texttt{emi: dto} definition and writes two files: the schema itself, and
a thin wrapper component around \href{https://rjsf-team.github.io/react-jsonschema-form/}{react-jsonschema-form}
(RJSF) that renders it:
\begin{lstlisting}
emi js:rjsf --path address.emi.yml --output ./out
# ./out/AddressDto.schema.json
# ./out/AddressDtoForm.tsx
\end{lstlisting}

\textbf{Embedded on the class} --- \texttt{--tags json-schema} on a full
module compile instead attaches every generated dto/action-request/
action-response class's schema directly to itself, as
\texttt{static JsonSchema = \{...\}} (off by default --- it's extra
bytes most consumers don't need):
\begin{lstlisting}
emi js --path module.emi.yml --output ./out --tags json-schema,typescript
\end{lstlisting}
Both read the exact same \texttt{BuildJSONSchema} tree walk
(\texttt{lib/formgen/jsonschema.go}), so the shape of the schema itself
never depends on which of the two you use --- only where it ends up.

\subsection{A minimal example}
\begin{lstlisting}
name: address
fields:
  - name: street
    type: string
    description: Street and number
  - name: unit
    type: string?
  - name: kind
    type: enum
    enum:
      - key: home
        description: Home address
      - key: work
  - name: manager
    type: one
    target: employee
\end{lstlisting}
compiles (\texttt{emi js:rjsf}) to:
\begin{lstlisting}
{
  "type": "object",
  "title": "AddressDto",
  "properties": {
    "street": { "type": "string", "title": "Street", "description": "Street and number" },
    "unit":   { "type": "string", "title": "Unit" },
    "kind":   { "type": "string", "title": "Kind", "oneOf": [
      { "const": "home", "title": "Home address" },
      { "const": "work", "title": "work" }
    ]},
    "manager": { "title": "Manager" }
  },
  "required": ["street", "kind", "manager"]
}
\end{lstlisting}

\begin{sidebar}{Two things worth noticing}
\texttt{unit} is missing from \texttt{required} --- the trailing
\texttt{?} is the only thing that removes a field from it, independent
of the field's own type or widget. And \texttt{manager} --- a
\texttt{one} relation, so genuinely unconstrained (\texttt{\{title:
"Manager"\}}, no \texttt{type}) --- is still \emph{in} \texttt{required},
because it isn't nullable: \texttt{required} tracks nullability alone,
not whether the schema actually constrains the value. RJSF only checks
the key is present, so an empty object still satisfies it.
\end{sidebar}

An \texttt{enum} option's human label comes from its \texttt{description}
(falling back to the raw key, as \texttt{work} does above); \texttt{one}/
\texttt{collection} relation fields and \texttt{any}/\texttt{complex}
fields both resolve to an unconstrained \texttt{\{\}} on purpose --- the
schema compiler was never given the target entity's own fields (a
relation) or has no fixed shape to begin with (\texttt{any}), so it
leaves the value alone rather than guessing, deliberately mirroring the
same ``leave it to a custom field on the renderer side'' choice
\S\ref{sec:js-complex} makes for a hand-written complex type.

\subsection{Using it in the front end}
The RJSF wrapper is a controlled component --- \texttt{value}/
\texttt{onChange}, the same convention every other generated form
component in Emi uses --- so a consuming app drops it in with zero
manual field wiring, validation included:
\begin{lstlisting}
import AddressDtoForm from "./out/AddressDtoForm";

function EditAddress() {
  const [value, setValue] = useState<Partial<AddressDto>>({});
  return <AddressDtoForm value={value} onChange={setValue} />;
}
\end{lstlisting}
Under the hood that wrapper is nothing more than the schema plus RJSF's
own \texttt{<Form>} and \texttt{@rjsf/validator-ajv8} --- swap either
import and the same generated \texttt{.schema.json} still works, since
nothing about it is RJSF-specific.

The embedded \texttt{static JsonSchema} form exists for the opposite
case: an app that isn't using RJSF at all still gets a real JSON Schema
for free on every generated class, to hand to a different schema-driven
renderer, validate a payload with a bare \texttt{ajv} instance before a
submit, or drive a headless dynamic-form engine --- all without a second
compile step or a hand-maintained schema that can drift from the dto.

\begin{pullquote}
Both paths are opt-in and additive: no code Emi already generates for a
dto changes shape because a schema exists alongside it, and skipping
both (the default) costs nothing --- no \texttt{formgen} import ends up
in the output at all.
\end{pullquote}
